\section{Related Work}
\label{slab:sec:related}

%\cancut{This section surveys work from both the databases and programming languages communities that is closely related to our paper.} 


% \head{Logical query rewrite}  Logical query rewrite optimizes a query by transforming a query into an equivalent one with better performance. Prior work mostly applies pattern-matching rules to perform query rewrite~\cite{wetune,LR,learning-to-rewrite,pirahesh1992extensible,chen2016memsql}. \barzan{overlap with intro.}

\head{Rule-based query rewriting} 
%Query rewriting is an important area of research  for  problem in query optimization, which transforms a query into an equivalent one with better performance. 
There is significant work on rule-based query rewriting~\cite{dintyala2020sqlcheck,lohman1988grammar,finance1991rule,wetune,pirahesh1992extensible,chen2016memsql,begoli2018apache}. 
% \jie{Maybe we can merge "Leveraging constraint for query optimization" here. }  \jie{Pattern-matching rules can also leverage integrity constraints for query optimization, for instance, using NOT NULL and key constraints to eliminate joins~\cite{join-elimination}, using NOT NULL constraint to optimize queries with disjunction~\cite{blog-disjunct}, using key constraints to remove unnecessary DISTINCT~\cite{habimana2015query}, and more~\cite{presto-qo,db2-qo}. }
The rules are either crafted by databases experts over decades~\cite{begoli2018apache}, which grow very slowly and cannot handle unanticipated patterns, or discovered by automated tools~\cite{wetune}, which can only handle a small subset of simple SQL queries.
In light of the recent interest~\cite{LR} in adopting deep learning in query optimization~\cite{yang2022balsa,krishnan2018learning,marcus2019neo}, there is also some work on using deep learning for query rewriting~\cite{LR}.
Given a SQL query and a set of rewrite rules, $\learnedrewrite$~\cite{LR} decides the order in which the rewrite rules should be applied using the Monte Carlo Tree method with learned cost models. However, both traditional and learning techniques are fundamentally limited by the incomplete nature of rules and pattern-matching; thus, they may miss valuable rewrite opportunities (see Section~\ref{slab:sec:motivating-example}).

\head{Leveraging constraints for query optimization} Prior work  explores semantic query rewriting by considering additional constraints. For instance, leveraging \texttt{NOT} \texttt{NULL} and key constraints to eliminate joins~\cite{join-elimination}, using \texttt{NOT} \texttt{NULL} constraint to optimize queries with disjunction~\cite{blog-disjunct}, considering key constraints to get rid of unnecessary \texttt{DISTINCT}~\cite{habimana2015query}, and more~\cite{presto-qo,db2-qo}. However, they mostly rely on pattern-matching rules to identify rewriting opportunities that leverage relatively simple constraints (e.g., foreign or primary keys), while \slabcity can flexibly use more complex forms of integrity constraints during query synthesis. 


\head{Checking query equivalence} 
There is a body of prior work on the automatic verification of query equivalence. Cosette~\cite{chu2017cosette} leverages proof assistants (e.g., Coq)
% \barzan{assistance or assistants or a proof assistant. which one did we mean? plz fix and remove this comment} \xinyu{Coq, proof assitant. the very PL stuff.}
to interactively construct mechanized proofs for equivalent query pairs or generate counterexamples for inequivalent pairs. EQUITAS~\cite{zhou2019automated,zhou2022spes}  executes queries symbolically to generate a logical formula that encodes the query semantics, and then uses an SMT solver to verify the query. It can be applied for a subset of SQL queries (SPJ and some outer join and aggregate). \slabcity is orthogonal to these techniques as it focuses on quickly finding examples that eliminate inequivalent query candidates. 
On top of that, $\slabcity$ leverages these techniques to check the equivalence of its candidate queries and provide formal guarantees of rewrite correctness. 

There is also work, from the software engineering and testing literature, that uses mutation-based testing techniques to identify common failure patterns of queries~\cite{chandra2015data,shah2011generating}. They use hard-coded constraint 
rules~\cite{shah2011generating} or specifications~\cite{chandra2015data} to generate test databases with the goal of killing as many query mutants as possible. These techniques consider a small query language that cannot express complex queries that arise from our workloads.
EvoSQL~\cite{castelein2018search} applies an evolutionary search algorithm that can test query correctness, in an offline manner, guided by the predicate coverage metric proposed by Tuya et al.~\cite{tuya2010full}. 
In contrast, the tester is $\toolname$ is guided by a set of common patterns that lead to incorrect rewrites. 
In other words, our approach is customized to our problem domain of query rewriting and can identify counterexamples more efficiently in an online fashion during the CEIGS-based synthesis process. 
%\tofix{\slabcity's goal is to identify incorrect synthesized queries rather than aiming for full predicate coverage of the input query.} 
%\tofix{prev sentence 1) is confusing and 2) makes the reviewer think oh gosh they only support very simple predicates compared to EvoSQL. plz clarify and rewrord what we are trying to say here}
RATest~\cite{miao2019explaining} uses a provenance-based algorithm to find a minimal database that can distinguish and explain incorrect queries; however, it is quite expensive and can only support a limited SQL language, which makes it not suitable for our domain. 

\begin{comment}


but this approach is expensive and can only support \tofix{a much simpler query language}.
\barzan{aren't they supporting SQL? if so we can say `a much smaller subset of SQL'}
\xinyu{Yuxuan may respond later but I think subset is correct.}

\yuxuan{On EvoSQL: I think EvoSQL does not fit in our scenario for two reasons. 1) EvoSQL is an offline approach, typically taking several minutes (stated in their paper) to generate databases for one input query. Thus, it's more time consuming than we expected. 2) The most important is that EvoSQL's target, the full predicate coverage, covers only a subset of what our synthesizer generates.

On RATest: RATest only support queries with the form of Select-Project-Join-Union-Difference, seems not runnable on queries with aggregation, groupby, window function, etc.}


\end{comment}



\head{Program synthesis for databases} 
Program synthesis techniques 
have been used in the area of databases. 
%, to optimize UDF queries, synthesize queries from input-output examples, or test database functionalities,
For example, 
Blitz~\cite{synthesis-spark,blitz2} synthesizes user-defined operators from Spark programs for better parallel query execution across operators.
Chestnut~\cite{chestnut} uses synthesis to find a new data layout as well as a query plan to execute on that layout. Unlike \slabcity, they focus on non-relational queries rather than optimizing core SQL queries.
\textsc{Sickle}~\cite{zhou2022synthesizing} and Scythe~\cite{wang2017synthesizing} synthesize analytical queries for end users given input-output examples, which lowers the burden on the user to remember the finer details of SQL.
% \jie{Sia is mentioned before; not sure if we mention it here. Sia~\cite{zhou2021sia} synthesizes valid predicates based on the given predicate and uses an SMT solver to verify the learned predicate. However, it limits its scope to only predicate synthesis and omits rewrite opportunities in other parts of queries while \slabcity does whole-query synthesis.} 
Program synthesis has also been used to test database applications~\cite{qex,tesma,syndb,xie:symbolic2}, where the goal is to automatically construct a more diverse set of test databases in order to improve the application code coverage.

%ta inja

\head{Data provenance} 
%Although not directly related to our work, 
We also briefly discuss how data provenance is related to (and different from) the concept of query dataflow in this paper, since they might look similar to each other. 
There is a large body of work on data provenance in various sub-fields of database systems, such as probabilistic databases~\cite{green2006models,suciu2011probabilistic}, view maintenance~\cite{buneman2002propagation}, and explanation of query results~\cite{green2007provenance,chapman2009not}. More recently, data provenance has also been used to guide the synthesis of Datalog programs~\cite{raghothaman2020provenance} and SQL queries~\cite{zhou2022synthesizing} from input-output examples. 
Data provenance is concerned with where every piece of data originates and how it is computed. 
In contrast, our notion of query dataflow is based on a coarser-grained, column-level flow of information within the input query.
As such, computing query dataflow is significantly more tractable than data provenance, yet sufficient enough to help guide \slabcity through its search.
